\begin{multicols}{2}
    \begin{exercice}{}{}
        Pour chacun des calculs suivants, donner l'opération prioritaire, simplifier l'expression en effectuant cette opération, et recommencer jusqu'à arriver au résultat final. Par exemple~:
        $$A = \underbrace{3 \times 4}_\text{prioritaire} + 5 = \underbrace{12 + 5}_\text{prioritaire} = 17$$
        \begin{enumeratebf}
            \item $A=7\times 12+5$
            \item $B=7\times (12+5)$
            \item $C=12,3-2,5+4$
            \item $D=8\times (8\times 2+4)\times 7$
        \end{enumeratebf}
    \end{exercice}



\begin{exercice}{}{}
    Effectuer les calculs suivants.
    \begin{enumeratebf}
        \item $A=(3+5)-(7-4)$
        \item $B=-(5+7)+(3-4)$
        \item $C=(3+5)-(7+4)$
        \item $D=(3-5)-(7+4)$
        \item $E=(3-7)-(4+5)$
    \end{enumeratebf}

\end{exercice}

\begin{exercice}{}{}
    Effectuer les calculs suivants.

    \begin{enumeratebf}
        \item $A=1,5\times 0,5+4\times (-2)$
        \item $B=(-1)\times (-2)\times (-3)\times (-4)$
        \item $C=8-8÷2$
        \item $D=14,2-8÷(-1)\times 2$
        \item $E=3\times (-5)-(-6)\times (-2)$
        \item $F=5\times (-2)+(-8)÷4$
    \end{enumeratebf}



\end{exercice}


\begin{exercice}{}{}
    Calculer les expressions suivantes, en respectant les priorités opératoires.
    \begin{enumeratebf}
        \item $A=3\times 4+5$
        \item $B=3\times (4+5)$
        \item $C=12,3-2,5+4$
        \item $D=8\times (8\times 2+4)\times 7$
    \end{enumeratebf}
\end{exercice}

\begin{exercice}{}{}
    Effectuer les calculs suivants.

    \begin{enumeratebf}
        \item $A=4\times (-3)+10$
        \item $B=(-5)-5\times 8+2$
        \item $C=6÷3-9\times 2$
        \item $D=6÷(3-9)\times 2$
    \end{enumeratebf}

\end{exercice}

\end{multicols}